% !TEX root = ../paper.tex

\section{Discussion}

\subsection{What UAV-TAlign-12K Reveals}

The main contribution of UAV-TAlign-12K is a captured dataset and protocol that expose a gap hidden
by pairwise-only reporting. Modern matchers provide abundant frame-level evidence on the official
split, yet the strict scene-level criterion retains only a subset of scenes. The benchmark
therefore separates two practically different capabilities: producing a pairwise homography and
deciding whether multi-frame evidence supports a trustworthy scene transform.

The protocol evidence supports this dataset-centered reliability framing. Retained and QA-filtered scenes show
structured differences in accepted ratio, severe outlier behavior, and baseline-relative QA
signals, indicating that the canonical criterion captures scene-level reliability beyond pairwise
matcher availability alone. This gives the evaluation a clear diagnostic role: it tests when
abundant pairwise evidence becomes a trustworthy scene-level transform. The canonical scene
criterion is therefore designed for conservative retention of scene transforms rather than for
maximizing unconditional scene acceptance.

UAV-TAlign serves as a reference implementation for this second evaluation unit. Its evidence
selection, robust consensus, and QA-aware verification demonstrate one reproducible way to convert
pairwise hypotheses into a retained scene product; the benchmark remains compatible with future
matchers and alternative scene-level aggregation strategies.

The multi-seed random-selection supplement further strengthens the design rationale for
deterministic evidence selection. On the accepted 8-scene subset, random selection spans 5/8 to
7/8 retained scenes across seeds, but the variation is concentrated in a few borderline scenes
rather than in large shifts of accepted-frame support. Deterministic-even selection therefore
provides a reproducible default operating policy, while robust scene consensus and QA-aware
verification define the retained-product decision layer.

\subsection{Scope, Limitations, and Intended Use}

UAV-TAlign-12K targets captured UAV scenes where a dominant planar or far-field transform is a
useful operating model. Its 15 scenes cover diverse illumination, thermal rendering, view, and
scene-family conditions, but they do not represent every flight altitude, sensor configuration, or
non-planar environment. The benchmark should therefore be interpreted as a focused resource for
cross-FOV homography alignment rather than a universal model of all UAV registration settings.

The label-free QA proxies enable scalable analysis without dense geometric annotations, but they
are scene-consistency diagnostics rather than substitutes for dense ground truth or downstream
task evaluation. The canonical gate is intentionally conservative, and the reliability score is
used only for ranking and coverage analysis. These boundaries make the protocol suitable for
comparing deployable scene-level reliability while leaving room for future annotations, geometric
models, and task-specific validation.
